\documentclass[letterpaper,12pt]{article}
\usepackage[T1]{fontenc}
\usepackage[utf8]{inputenc}
\usepackage[spanish]{babel}
\usepackage{framed}
\usepackage{amsfonts,setspace}
\usepackage{amsmath}
\usepackage{amssymb, amsmath, amsthm}
\usepackage{comment}
\usepackage[most]{tcolorbox}
\usepackage{amssymb}
\usepackage{dsfont}
\usepackage{anysize}
\usepackage{pdfpages}
\usepackage{multicol}
\usepackage{enumerate}
\usepackage{graphicx}
\usepackage{float}
\usepackage{tikz}  
\usepackage[dvipsnames]{xcolor}
\usepackage[left=1.5cm,top=0.9cm,right=1.5cm, bottom=1.4cm]{geometry}
\usepackage{fancyhdr}
\pagestyle{fancy}
\definecolor{blue(pigment)}{rgb}{0.2, 0.2, 0.6}
% Page
\oddsidemargin -0.4in
\textwidth 7in
\topmargin -0.6in
\textheight 9.1in
\parindent 0em
\parskip 1ex

%Teoremas, Lemas, etc.
\theoremstyle{plain}
\newtheorem{teo}{Teorema}
\newtheorem{lem}{Lemma}
\newtheorem{prop}{Proposici\'on}
\newtheorem{cor}{Corolario}
\theoremstyle{definition}
\newtheorem{defi}{Definici\'on}
\newtheorem{eje}{Ejemplo}
\newtheorem{obs}{Observaci\'on}
% fin macros

% Comandos más usados
% Conjuntos
\newcommand{\N}{\mathbb{N}}
\newcommand{\Z}{\mathbb{Z}}
\newcommand{\Q}{\mathbb{Q}}
\newcommand{\R}{\mathbb{R}}
\newcommand{\C}{\mathbb{C}}
\newcommand{\RR}{\mathbb{R} \times \mathbb{R}}
% Flechas
\newcommand{\ssi}{\Longleftrightarrow}
\newcommand{\imp}{\Longrightarrow}
\newcommand{\pmi}{\Longleftarrow}
% Trigonométricas
\newcommand{\sen}{\text{sen}}

\usepackage{versions}
\includeversion{solution}

\let\epsilon\varepsilon

\begin{document}
\def\Pauta{1}
%%%%%ESCRIBIR 1 si quieren ver la Prueba con Pauta%%%%
%%%%%ESCRIBIR 0 si quieren ver sólo la Ayudantía%%%%
	
	\fancyhead[L]{\itshape{Escuela de Ingeniería}}
	\fancyhead[R]{\itshape{Universidad de O'Higgins}}
	
	\begin{minipage}{11.5 cm}
         \vspace{-6mm}
		\begin{flushleft}
			\vspace{1mm}
			\textbf{ING1002-2 Cálculo Diferencial e Integral}\\
			\textbf{Profesor:}  Gonzalo Flores G.  \\
			\textbf{Ayudante:}  Juan I. Riquelme \& Cristian Acevedo R. \\
		\end{flushleft}
	\end{minipage}
	
	\begin{picture}(2,3)
		\put(430,9.5){\includegraphics[scale=0.22]{Imagenes/UOH.png}}
	\end{picture}
	\vspace{-6mm}
	\begin{center}
		\LARGE \bf{Ayudantía 13: Integrales Impropias} %Teorema Fundamental del Cálculo
	\end{center}

\begin{enumerate}[\textbf{P\arabic{enumi}.}]
\item Calcule las siguientes integrales impropias y determine si son o no convergentes:
    \begin{enumerate}
        \item $\displaystyle \int_0^{\infty} \frac{x^2}{\sqrt{1 + x^3}} \, dx$

\color{blue}
\textbf{Sol:} Cambio de variable
\[
u=1+x^3 \;\Rightarrow\; du=3x^2\,dx.
\]
\[
\int_0^{\infty}\frac{x^2}{\sqrt{1+x^3}}\,dx
=\frac{1}{3}\int_{1}^{\infty} u^{-1/2}\,du.
\]
\[
\frac{1}{3}\int u^{-1/2}du
=\frac{1}{3}\cdot 2u^{1/2}
=\frac{2}{3}\sqrt{u}.
\]

\[
\lim_{b\to\infty}\frac{2}{3}\big(\sqrt{b}-1\big)=\infty.
\]
tenemos que la integral es \textbf{divergente}.
\color{black}
        
        \item $\displaystyle \int_0^1 \ln(x)dx$

\color{blue}
\textbf{Sol:} $\displaystyle \int_{0}^{1} \ln(x) \, dx \quad \leadsto$ es de tipo 2 ya que $\displaystyle \lim_{x \to 0} \ln(x) \to -\infty$
Calculemos directo:
\begin{align*}
\int_{0}^{1} \ln(x) \, dx &= \lim_{a \to 0} \int_{a}^{1} \ln(x) \, dx \\
&= \lim_{a \to 0} (x \ln(x) - x) \bigg|_{a}^{1} \\
&= \lim_{a \to 0} \left[ (1 \ln(1) - 1) - (a \ln(a) - a) \right] \\
&= -1
\end{align*}
para el límite auxiliar:
\[
\begin{aligned}
\lim_{a \to 0} a \ln(a) &= \lim_{a \to 0} \frac{\ln(a)}{\frac{1}{a}} \\
&\stackrel{L'H}{=} \lim_{a \to 0} \frac{\frac{1}{a}}{-\frac{1}{a^2}} \\
&= \lim_{a \to 0} (-a) = 0
\end{aligned}
\]
Por lo tanto, el resultado es $-1$. \textbf{¡Converge!}
\color{black}
        
        \item $\displaystyle \int_0^{\infty} \frac{dx}{x(x+\sqrt{x})}$

\color{blue}
\textbf{Sol:} $\displaystyle \int_{0}^{\infty} \frac{dx}{x(x+\sqrt{x})} = \int_{0}^{a} \frac{dx}{x(x+\sqrt{x})} + \int_{a}^{\infty} \frac{dx}{x(x+\sqrt{x})} \quad \text{con } a > 0$.  Si $a = 1$, estudiemos $(I)$:
\[ \underbrace{\int_{0}^{1} \frac{dx}{x(x+\sqrt{x})}}_{(I)} + \underbrace{\int_{1}^{\infty} \frac{dx}{x(x+\sqrt{x})}}_{(J)} \]
Estudiemos (I): Para $x \in [0, 1]$, se cumple que $x \leq \sqrt{x}$.
\[ \implies \frac{1}{x(x+\sqrt{x})} \geq \frac{1}{x(\sqrt{x}+\sqrt{x})} = \frac{1}{2x \cdot x^{1/2}} = \frac{1}{2} x^{-3/2} \]
Entonces:
\[ \implies \int_{0}^{1} \frac{1}{x(x+\sqrt{x})} \, dx \geq \int_{0}^{1} \frac{1}{2} \frac{1}{x^{3/2}} \, dx \]
Como $\frac{3}{2} > 1$ (en el límite inferior), la integral de la derecha \textbf{diverge}.

$\therefore \displaystyle \int_{0}^{1} \frac{1}{x(x+\sqrt{x})} \, dx$ diverge $\displaystyle \implies \int_{0}^{\infty} \frac{1}{x(x+\sqrt{x})} \, dx$ también.
\color{black}
        
        \item $\displaystyle \int_{0}^\infty \frac{\arctan(x)}{x^2} \, dx$

\color{blue}
\textbf{Sol:} $\displaystyle \int_{0}^{\infty} \frac{\arctan(x)}{x^2} \, dx \ \text{es} \ \text{mixta}$
\noindent $\displaystyle = \underbrace{\int_{0}^{1} \frac{\arctan(x)}{x^2} \, dx}_{\text{2\textsuperscript{a} especie}} + \underbrace{\int_{1}^{\infty} \frac{\arctan(x)}{x^2} \, dx}_{\text{1\textsuperscript{a} especie}}$
\quad 

\begin{minipage}[t]{0.4\textwidth}
\textit{En cociente el límite es $x \to 0$ pues en 0 explota la función.}
\end{minipage}
\noindent Veamos esta (la de 2\textsuperscript{a} especie),

usemos \textbf{criterio del cociente}:
\[
\lim_{x \to 0} \frac{\left( \frac{\arctan(x)}{x^2} \right)}{\frac{1}{x}} = \lim_{x \to 0} \frac{\arctan(x)}{x} \stackrel{L'H}{=} \lim_{x \to 0} \frac{\frac{1}{x^2+1}}{1} = 1 \neq 0
\]
\noindent Por lo tanto $\displaystyle \int_{0}^{1} \frac{\arctan(x)}{x^2} \, dx$ \textbf{diverge} ya que $\displaystyle \int_{0}^{1} \frac{1}{x} \, dx$ \textbf{diverge}. $\implies \displaystyle \int_{0}^{\infty} \frac{\arctan(x)}{x^2} \, dx$ también \textbf{diverge}.
\color{black}
        
    \end{enumerate}

\item Para el siguiente conjunto $S = \{(x, y) : x \geq 0, 0 \leq y \leq \lambda e^{-\lambda x}, \lambda > 0\}$, trace su región, y determine si esta es finita o no.

\color{blue}
\textbf{Sol:} El conjunto dado es
\[
S = \{(x,y) : x \ge 0, \; 0 \le y \le \lambda e^{-\lambda x}, \; \lambda>0\}.
\]
Para cada $x \ge 0$, $y$ varía entre $0$ y $y=\lambda e^{-\lambda x}$. La curva superior es una función exponencial decreciente que parte de $y=\lambda$ en $x=0$ y tiende a $0$ cuando $x\to\infty$. La curva inferior es $y=0$ (eje $OX$).
La región $S$ es
\[
\begin{tikzpicture}[scale=1.0]
\draw[->] (-0.2,0) -- (6,0) node[right] {$x$};
\draw[->] (0,-0.2) -- (0,3) node[above] {$y$};
\draw[thick,domain=0:5,smooth,samples=100] plot (\x,{2*exp(-0.5*\x)}) node[right] {$y=\lambda e^{-\lambda x}$};
\draw[fill=blue!20,domain=0:5,smooth,samples=100] 
  plot (\x,{2*exp(-0.5*\x)}) -- (5,0) -- (0,0) -- cycle;
\end{tikzpicture}
\]
calculamos su área:
\[
A = \int_0^{\infty} \lambda e^{-\lambda x}\,dx = \left[-e^{-\lambda x}\right]_0^{\infty} = 1.
\]
Conclusión: la región es \textbf{finita} y tiene área $1$.
\color{black}

\item \textbf{[Propuesto].-} Calcule, si es que existe, el área de bajo de la curva
$$f(x)=\frac{x}{(1+x^2)^{3/2}}$$
En el primer cuadrante.

\color{blue}
\textbf{Sol:} Nos piden $A = \displaystyle \int_{0}^{\infty} f(x) \, dx = \int_{0}^{\infty} \frac{x}{(1+x^2)^{3/2}} \, dx$
Vayamos directo a calcular. Obs: es de \textbf{1era especie}.
\begin{align*}
\int_{0}^{\infty} \frac{x}{(1+x^2)^{3/2}} \, dx &= \lim_{b \to \infty} \int_{0}^{b} \frac{x}{(1+x^2)^{3/2}} \, dx \quad \leadsto \text{hacer desaparecer la raíz} \\
&\text{Cambio de variable:} \\
&u^2 = 1+x^2 \implies 2u \, du = 2x \, dx \implies u \, du = x \, dx \\
&\text{Límites:} \\
&\text{Si } x \to \infty \implies u \to \infty \\
&\text{Si } x = 0 \implies u = 1 \\
\\
&= \lim_{b \to \infty} \int_{1}^{b} \frac{u \, du}{(u^2)^{3/2}} \\
&= \lim_{b \to \infty} \int_{1}^{b} \frac{u \, du}{u^3} \\
&= \lim_{b \to \infty} \int_{1}^{b} \frac{du}{u^2} \quad \text{es del tipo } \frac{1}{x^p} \text{ con } p=2 > 1, \text{ por lo que converge.} \\
&= \lim_{b \to \infty} \left( \frac{u^{-2+1}}{-2+1} \right) \bigg|_{1}^{b} \\
&= \lim_{b \to \infty} \left( -\frac{1}{u} \right) \bigg|_{1}^{b} = \lim_{b \to \infty} \left[ -\frac{1}{b} - \left( -\frac{1}{1} \right) \right] \\
&= 1
\end{align*}
\color{black}
\end{enumerate}

\end{document}